\subsection{Local rings of equal characteristics}

Let A be a ring. Recall (A, V, p. 2) that the characteristic of A is defined when A contains a subfield. It is equal to 0 if and only if A contains a subfield isomorphic to Q, and equal to a prime number p if and only if one has \(p1_{A} = 0\) . If the characteristic of A is defined, and if \(f: A \to B\) is a nonzero homomorphism of rings, the characteristic of B is defined and is equal to that of A.

Let A be a local ring, with maximal ideal m, and residue field k.

\begin{enumerate}
\def\labelenumi{\alph{enumi})}
\item
  Suppose that \(k\) has characteristic 0. Then A contains a field and the characteristic of A is equal to 0. Indeed, the canonical homomorphism from Z to A is injective, and for every nonzero integer \(n\) , \(n1_{\mathbf{A}}\) is invertible in A, because it does not belong to m.
\item
  Suppose that \(k\) has characteristic \(p \neq 0\) . Then A contains a field if and only if \(p1_{A} = 0\) . In this case the characteristic of A is equal to \(p\) .
\end{enumerate}

Suppose that A is an integral local ring, with field of fractions K and residue field k.

\(a^{\prime})\) The ring A contains a subfield if and only if the characteristics of \(k\) and K are equal. In this case, the characteristic of A is equal to that of \(k\) and of K, and one says that A is a local ring of equal characteristics.

\(b^{\prime})\) Suppose that the fields \(k\) and \(\mathbf{K}\) do not have the same characteristic. Then there exists a prime number \(p\) such that \(k\) is of characteristic \(p\) . Since one has \(q1_{\mathbf{A}} \neq 0\) for every prime number \(q \neq p\) , the field \(\mathbf{K}\) is of characteristic 0. One then says that A is a local ring of unequal characteristics.

